% cuadro C4_4
\begin{table}[htbp]
\centering
\scriptsize
\caption{Gastos obligatorios con y sin pensiones, y el agregado que su diferencia adjudica}
\label{cua:C4_4}
\begin{tabular}{p{5.2cm}rrrrrr}
\toprule
\textbf{Concepto} & \textbf{Proy. 2025} & \textbf{Proy. 2026} & \textbf{Dif. nom.} & \textbf{Dif. real} & \textbf{Var. real \%} & \textbf{\% PIB} \\
\midrule
Gastos obligatorios sin pensiones & 6 046 446,8 & 6 702 281,0 & 655 834,2 & 365 604,8 & 5,77 & 17,31 \\
Gastos obligatorios con pensiones y jubilaciones & 7 684 111,9 & 8 406 440,2 & 722 328,3 & 353 490,9 & 4,39 & 21,71 \\
Pensiones y jubilaciones (diferencia) & 1 637 665,1 & 1 704 159,2 & 66 494,1 & -12 113,8 & -0,71 & 4,40 \\
\bottomrule
\end{tabular}
\\[2pt]\begin{minipage}{0.92\textwidth}\scriptsize La diferencia entre los dos primeros renglones es el agregado de pensiones y jubilaciones de la clasificación económica. Es una fuente independiente que adjudica el perímetro y no una elección metodológica. Fuente: Anexo 3 del proyecto de decreto de 2026 y de 2025. Elaborado por el ITED.\end{minipage}
\end{table}
